\documentclass[12pt]{article}
%^ Start of document
\usepackage{times}  %Makes text TimesNewRoman
\usepackage{setspace} %Allows you to set single or double space 
\usepackage{biblatex} %Imports biblatex package
\usepackage{graphicx} % Required for inserting images
\usepackage{authblk}
\usepackage[colorlinks=true, urlcolor=blue, linkcolor=blue]{hyperref}
\usepackage{subfig}
\usepackage{float}
\usepackage{tabularx}
\usepackage{multirow}
\usepackage[paperheight=11in,
   paperwidth=8.5in,
   top=20mm,
   bottom=20mm,
   left=40mm,
   right=40mm]{geometry}
\usepackage{moresize}
\usepackage{tikz}
\usepackage{xcolor}
\usepackage{physics}
\usepackage{amssymb}
\usepackage{mathrsfs}
\usepackage{amsmath}
\usepackage[document]{ragged2e}
\usepackage{comment}
\usepackage{xcolor}
\addbibresource{Qhw.bib}
\begin{document}
\singlespacing  %Sets single spacing for whole document
\title{The Quest for Unification in Modern Physics}

\author[1]{Zachary Tullier}
\affil[1]{Student\\Physics Department\\ University of Houston - Clear Lake \\Houston, Texas\\U.S}
\date{}
\maketitle


\section*{Abstract}
This paper summarizes a lecture by Dr. Qaisar Shafi on the theme of unification in physics---both historical and theoretical. Dr. Shafi traces the evolution of fundamental physics from classical electromagnetism through special and general relativity, to quantum field theory and the Standard Model. He outlines the theoretical framework and discoveries that unify fundamental forces and particles, culminating in a contemporary understanding of high-energy physics and cosmology. The lecture also explores unresolved questions such as the nature of dark matter and dark energy, the role of quantum fluctuations in structure formation, and the elusive magnetic monopole.

\section*{Introduction: A Journey Toward Unification}
Dr. Shafi’s presentation is a pedagogically rich journey through the major milestones of modern physics, emphasizing the historical and theoretical development of \textbf{unification}---the idea that seemingly distinct physical forces and particles may be manifestations of a single underlying theory. He begins by establishing the conceptual groundwork, introducing natural units and fundamental constants such as the Planck scale, and setting the stage for quantum gravity and cosmology.

\section*{Electromagnetism and Special Relativity}
Unification begins with James Clerk Maxwell’s 19th-century synthesis of electricity and magnetism into a single theory of \textbf{electromagnetism}, expressed in a set of elegant equations. This unification led to the prediction of electromagnetic waves and established a key principle: the \textbf{Lorentz invariance} of physical laws. This concept underpins Einstein’s 1905 theory of \textbf{special relativity}, which postulates the constancy of the speed of light for all observers and introduced the notion of four-dimensional spacetime.

\section*{General Relativity and Dynamical Spacetime}
Einstein’s \textbf{general theory of relativity} (1915) generalized special relativity to include gravity as a curvature of spacetime, governed by the famous field equation:
\[
G_{\mu\nu} = T_{\mu\nu}
\]
Here, gravity is no longer a force in the Newtonian sense, but a geometric property of spacetime influenced by energy and momentum. General relativity set the foundation for modern cosmology and the understanding of massive structures like black holes and expanding universes.

\section*{Quantum Mechanics and the Dirac Equation}
As the quantum revolution emerged in the early 20th century, physicists developed models to describe the probabilistic nature of atomic and subatomic phenomena. Dr. Shafi details \textbf{Dirac’s synthesis of quantum mechanics and special relativity}, culminating in the \textbf{Dirac equation}, which not only predicted the existence of antimatter (e.g., the positron) but also provided the first relativistic quantum field theory for spin-$\frac{1}{2}$ particles.

\section*{Quantum Field Theory and the Standard Model}
Post-WWII, physicists formalized \textbf{quantum field theory (QFT)} to describe particle interactions via field excitations. QFT forms the basis of the \textbf{Standard Model (SM)}, which unifies three of the four fundamental forces:
\begin{itemize}
    \item \textbf{Strong interaction} (quantum chromodynamics, QCD)
    \item \textbf{Electromagnetic interaction}
    \item \textbf{Weak interaction}
\end{itemize}

These forces are governed by the gauge groups:
\begin{itemize}
    \item \( SU(3) \) for QCD (gluons),
    \item \( SU(2) \times U(1) \) for the electroweak interaction (W, Z bosons, and photon).
\end{itemize}

The \textbf{Higgs boson}, predicted by the electroweak theory and discovered at CERN in 2012, provides mass to gauge bosons via spontaneous symmetry breaking.

\section*{Structure of Matter: Quarks and Leptons}
Dr. Shafi reviews the composition of matter in terms of \textbf{three families} of quarks and leptons, each with increasing mass:
\begin{itemize}
    \item Quarks come in six flavors and three color charges.
    \item Leptons include electrons, muons, taus, and their corresponding neutrinos.
\end{itemize}

Protons and neutrons are bound states of quarks held together by gluons, confined by color charge. The non-observability of free quarks is a feature of \textbf{color confinement} in QCD.

\section*{Cosmology and the Early Universe}
Using the standard Big Bang model, Dr. Shafi discusses the universe’s thermal history. By combining general relativity with the Standard Model, physicists construct a consistent picture of the early universe from \( 10^{-12} \) to \( 10^{-35} \) seconds after the Big Bang. The \textbf{inflationary epoch}, driven by a scalar inflaton field (possibly related to the Higgs), explains homogeneity, isotropy, and the origin of \textbf{quantum fluctuations} that seed galaxies.

\section*{Open Questions: Dark Matter, Dark Energy, and Monopoles}
The lecture concludes with several profound questions still facing physics:
\begin{itemize}
    \item \textbf{Dark matter} constitutes approximately 25\% of the universe but remains undetected in particle physics.
    \item \textbf{Dark energy}, modeled as a cosmological constant \( \Lambda \), drives the universe’s current accelerated expansion and comprises about 70\% of the cosmic energy budget.
    \item \textbf{Magnetic monopoles}, predicted by grand unified theories and topological arguments, have yet to be observed, though experimental searches continue at CERN and in cosmic ray detectors.
\end{itemize}

Dr. Shafi also touches on the group-theoretic description of gravity and contrasts its complexity with gauge theories, explaining why gravitation is not simply incorporated into existing gauge frameworks---yet might find a natural place in string theory.

\section*{Conclusion}
Dr. Shafi’s talk masterfully weaves together the historical development, theoretical sophistication, and experimental confirmation of unification in modern physics. From Maxwell to Higgs, from Dirac to cosmological inflation, he demonstrates how humanity’s quest to understand nature reduces complexity into symmetry, and how much remains to be uncovered---especially in the realm of dark matter, quantum gravity, and early-universe physics.

\end{document}

\end{document}
